%! Constructing a Confidence Interval for a Population Proportion — Unit 6, Lesson 6.2 — Homework (Answer Key)
\documentclass[10pt]{article}
\usepackage{apstats-article}
\usepackage{apstats-key}   % pulls in -boxes; do NOT also load -boxes
\newcommand{\TallMath}[1]{$\displaystyle #1\rule[-1.4em]{0pt}{3.2em}$}

\begin{document}

\pageheader{Unit 6, Lesson 6.2}{Homework --- Answer Key}
\namedateperiod

\vspace{0.15in}

\textit{For each problem, use the four-step procedure: State, Plan (verify conditions), Do (compute interval), Conclude (interpret).}

\begin{enumerate}

  \item A random sample of 400 registered voters in a city finds that 228 support a proposed
        park expansion. Construct a 95\% confidence interval for the true proportion of all
        registered voters in the city who support the expansion.

        \textbf{State:} \ansline{$p$ = true proportion of all registered voters in the city who support the park expansion}
        \textbf{Plan:}
        \begin{itemize}
          \item Random: \ans{SRS stated --- met}
          \item Large Counts: \ans{$400(0.57)=228\ge10$; $400(0.43)=172\ge10$ --- met}
          \item 10\% rule: \ans{$400 \le 0.10\times$(all registered voters in the city) --- met}
        \end{itemize}
        \textbf{Do:} $\hat p = $ \ans{$228/400=0.57$}
        \[
          CI = \ans{0.57} \pm \ans{1.960}\sqrt{\frac{\ans{0.57(0.43)}}{\ans{400}}}
             = (\ans{0.521},\; \ans{0.619})
        \]
        \textbf{Conclude:} \ansline{We are 95\% confident the interval (0.521, 0.619) captures the true proportion of all registered voters in the city who support the park expansion.}

  \item A wildlife biologist tags and releases 600 randomly selected fish from a lake.
        Later, a random sample of 80 fish from the lake is taken, and 12 are tagged.
        Construct a 99\% confidence interval for the true proportion of tagged fish in the lake.

        \textbf{State:} \ansline{$p$ = true proportion of fish in the lake that are tagged}
        \textbf{Plan:}
        \begin{itemize}
          \item Random: \ans{SRS stated --- met}
          \item Large Counts: \ans{$80(0.15)=12\ge10$; $80(0.85)=68\ge10$ --- met}
          \item 10\% rule: \ans{$80 \le 0.10\times$(all fish in the lake) --- met}
        \end{itemize}
        \textbf{Do:} $\hat p = 12/80 = 0.15$
        \[
          CI = 0.15 \pm 2.576\sqrt{\frac{0.15(0.85)}{80}}
             = (\ans{0.047},\; \ans{0.253})
        \]
        \textbf{Conclude:} \ansline{We are 99\% confident the interval (0.047, 0.253) captures the true proportion of fish in the lake that are tagged.}

  \item A school district wants to estimate the proportion of families who would use
        an after-school tutoring program. An administrator wants a 95\% confidence interval
        with a margin of error of at most 0.05. A pilot survey suggests $\hat p \approx 0.40$.
        \begin{enumerate}[label=\alph*.]
          \item Find the minimum required sample size using the pilot estimate.
                \[n \ge \left(\frac{z^*}{m}\right)^2 \hat p(1-\hat p) = \left(\frac{1.960}{0.05}\right)^2(0.40)(0.60) = 1536.64\]
                Minimum $n = $ \ans{1537}

          \item How would your answer change if you used $\hat p = 0.50$ instead? Would it be
                larger or smaller?
                \ansline{Larger: $n \ge (1.960/0.05)^2(0.50)(0.50) = 1536.64 \to 1537$. Wait --- here $0.5(0.5)=0.25 > 0.40(0.60)=0.24$, so $n$ increases to 1537 (same rounding here, but $n$ is larger with $\hat p=0.5$ in general because $p(1-p)$ is maximized at 0.5).}
        \end{enumerate}

\end{enumerate}

\begin{reflectionbox}
\textbf{Preview:} In Lesson 6.3 you will focus on \emph{interpreting} confidence intervals ---
what ``95\% confident'' really means for the method, and how to use an interval to evaluate a
researcher's claim about a population proportion. Think: if your CI does \emph{not} contain a
claimed value, what conclusion can you draw?
\end{reflectionbox}

\end{document}
